%DS0



\section{ĐỀ SỐ \thesection }
\setcounter{bt}{0}
%%Câu 1.
\begin{bt} \textbf{}
	
	\begin{listEX}
		\item Thực hiện phép tính: $[53-(-12) \cdot 3 \cdot(-10)]$;
		\item Bỏ ngoặc rồi tính: $(177+222-185)-(67+22-75)$;
		\item Tìm số nguyên $x$ biết: $(x+5)(15x-30)=0$.
	\end{listEX}
	
\loigiai{a) $-307$ \hfill b) $200$ \hfill c) $x=-5, x=2$ 	}
\end{bt}


%%Câu 2.
\begin{bt} \textbf{ }
	\begin{listEX}
		\item  
		Tìm các số tự nhiên $x$ biết $42 \, \vdots \,  x$, $105 \, \vdots \, x$, $84 \, \vdots \, x$ và $x$ là bội của 3 . 
		\item 
		Số học sinh khối 6 của một trường trung học cơ sở trong khoảng từ $350$ đến $450$ em. Cô tổng phụ trách thấy rằng khi học sinh khối 6 xếp hàng $9$, hàng $10$, hàng $12$ đều thừa ra $3$ học sinh. Hỏi khối 6 của trường đó có bao nhiêu học sinh?
		
	\end{listEX}
	
	\loigiai{a) $x\in\{3;21 \}$ \hspace{2.5 cm} b) $363$ 	}
\end{bt}



%%Câu 3.
\begin{bt}
	Nhiệt độ trong một tủ cấp đông khi chưa bật là $23^{\circ} \mathrm{C}$. Khi bật tủ cấp đông, nhiệt độ bên trong tủ giảm $2^{\circ} \mathrm{C}$ mỗi phút. Hỏi phải mất bao nhiêu phút để nhiệt độ trong tủ đông đạt $-11^{\circ} \mathrm{C}$
	\loigiai{$17$ phút.}
\end{bt}


%%Câu 4.
\begin{bt}
	\immini[thm]{
		Trên một mảnh đất hình chữ nhật có chiều dài $12$ m, chiều rộng $10$ m người ta chia khu để trồng hoa, trồng cỏ: hoa sẽ được trồng ở khu vực hình bình hành $AMCN$, cỏ sẽ được trồng ở phần đất còn lại như hình vẽ. Tiền công để trả cho mỗi mét vuông trồng hoa là $50\,000$ đồng, trồng cỏ là $40\,000$ đồng.
		\begin{listEX}
			\item Tính diện tích mảnh đất hình chữ nhật, diện tích mảnh đất trồng hoa nêu trên.
			\item Tính số tiền công cần chi trả để trồng hoa và cỏ.
		\end{listEX}
	}{
		\begin{tikzpicture}[scale=1,line width=0.8pt,line join=round,>=stealth]
			\def\a{4.8} \def\b{4} 
			\path (0,0) coordinate (B)
			(0,\b) coordinate (A)
			(\a,\b) coordinate (D)
			(\a,0) coordinate (C)
			($(A)!0.5!(D)$) coordinate (N)
			($(B)!0.5!(C)$) coordinate (M)
			;
			\draw 
			(A)--(N)node[midway,above]{6 m}
			(D)--(N)node[midway,above]{6 m}
			(B)--(M)node[midway,below]{6 m}
			(M)--(C)node[midway,below]{6 m}
			(A)--(B)node[midway,left]{10 m}
			(A)--(M) (N)--(C)--(D)
			;
			\foreach \x/ \goc in {A/135,B/-135,C/-45,D/45,N/90,M/-90}  \fill (\x) circle (0pt)        
			($(\x)+(\goc:3.5mm)$) node {$\x$};
		\end{tikzpicture}
	}
	
	
	\loigiai{a) Diện tích mảnh đất là $120 $m$^2$, diện tích trồng hoa là $60 $m$^2$ \hspace{ 1 cm}
		b) $5\, 400\, 000$ }
\end{bt}


%%Câu 5.
\begin{bt}
	\immini[thm]{Cho biểu đồ cột kép biểu diễn số lượng học sinh yêu thích các môn thể thao của lớp 6A5
		\begin{listEX}
			\item Hãy cho biết môn thể thao nào có số lượng học sinh nữ và nam thích bằng nhau 
			\item Hãy cho biết các môn thể thao nào có số lượng học sinh nữ thích ít hơn số lượng học sinh nam và ít hơn bao nhiêu bạn.
		\end{listEX}
	}
	{\begin{tikzpicture}[scale=1,font=\small,line join=round,>=stealth]
			\def\hoanh{8} \def\tung{4} \def\sc{3.5}
			\def\mau{mauChuDao};
			\draw[->] (0,0)node[left]{0}--(\hoanh,0);
			\draw[->] (0,0)--(0,\tung+1);
			\foreach \y in {2,4,...,14}{\draw (-0.1,\y/\sc)node[left]{$\y$}--++(0:0.2);
				\draw[dashed,gray] (0,\y/\sc)--(\hoanh-0.2,\y/\sc);}
			\foreach \x/\n/\p in{1/6/13,3/10/10,5/12/10,7/14/5}{
				\draw[line width=5mm,\mau] (\x-0.3,0)--++(0,\n/\sc);
				\draw[line width=5mm,cyan] (\x+0.3,0)--++(0,\p/\sc);}
			\foreach \x/\p in {1/Bóng đá,3/Bóng bàn,5/Bóng rổ,7/Cầu lông}{\node at (\x,-0.4){\footnotesize \p};}        
			\fill[\mau] (1.5,\tung+0.8) rectangle (1.8,\tung +0.5) node [black,shift={(10:0.9)}]{Nữ};
			\fill[cyan] (4.5,\tung+0.8) rectangle (4.8,\tung +0.5) node [black,shift={(10:0.9)}]{Nam};
		\end{tikzpicture}
	}
	
	
	\loigiai{
	a) Bóng bàn \hfill b) Bóng đá (ít hơn $7$)
	 }
\end{bt}

%%Câu 6.
\begin{bt} So sánh
	$2^{777}$ và $7^{222}$.
	
	\loigiai{ $2^{777} > 7^{222}$. }
\end{bt}